% status: ZERO
% Printed Chapter 36 / stable source ch19. Rules: AUTHORING.md.
\chapter{Scheduling Under Pressure: Preemption, Fairness and SLO-Aware Admission}\label{ch:scheduling}

\begin{ChapterIntent}
A serving engine is easy to schedule while every runnable request fits. The real
scheduler appears when token budget, KV capacity and latency objectives cannot all be
satisfied. This chapter separates three decisions that are often conflated:
\emph{admission} decides whether new work may enter; \emph{scheduling} divides each
iteration among admitted work; and \emph{reclamation/preemption} takes resources back
after optimistic assumptions fail. We derive overload behavior, compare recomputation
with swapping, construct fair and priority-aware policies, and optimize for goodput
rather than raw throughput.
\end{ChapterIntent}

\OpeningSection{One exhausted cache, four simultaneous failures}

The edition's M1 measurement record contains an unusually sharp overload experiment.
llama.cpp's server was configured with one shared KV cache of 4,096 cells and four
request slots. Four requests arrived together, each with roughly a 200-token prompt
and permission to generate 1,200 tokens. Their eventual demand was about 5,600 cells.
After each had generated 824 tokens, the next step could not find enough KV space.
The server reduced the batch repeatedly, reached a one-token attempt, and then failed
all four active requests together with a context-size error. They failed at the same
pressure point, about 73.5 seconds after arrival (measured; record
\texttt{research/measurements/2026-09-24-m1-part3.md}). With 8,192 cells, the four
requests completed.

The scarce resource was real. The correlated failure was policy.

A scheduler with preemption could have selected one victim, reclaimed its state, and
let the others progress. A stricter admission controller could have delayed or
rejected work before all four accumulated expensive state. A reservation-based system
could have guaranteed completion but admitted fewer requests. None removes scarcity.
They decide who bears it, when, and how much useful work is destroyed.

\Foundations{Three control planes, three different questions}{
Admission asks ``may this request enter the resource domain?'' Scheduling asks ``which
admitted work executes in the next iteration?'' Reclamation asks ``which resident
state may be removed when earlier assumptions no longer fit?'' Chapter~\ref{ch:continuous-batching}
defines the iteration, Chapter~\ref{ch:paged-kv} the physical KV blocks,
Chapter~\ref{ch:chunked-prefill} the per-iteration work budget, and
Chapter~\ref{ch:prefix-caching} the reusable-state contract. A queue does not create
capacity. Preemption does not stabilize a persistently overloaded arrival stream.
Priority does not by itself provide fairness.
}

\NapkinSection{Napkin math: why 95 percent can feel saturated}

For a stable system, Little's law states
\begin{equation}
N=\lambda T,
\end{equation}
where $N$ is the average number of requests in the system, $\lambda$ the completed
arrival rate in requests/s, and $T$ mean time in system in seconds. It is an accounting
identity under its stability assumptions.

A deliberately simplified M/M/1 queue gives
\begin{equation}
T=\frac{1/\mu}{1-\rho},\qquad \rho=\frac{\lambda}{\mu},
\label{eq:sched-mm1}
\end{equation}
with service rate $\mu$ and utilization $\rho$. At $\rho=0.8$, mean time is five
service times; at $0.95$, twenty; at $0.98$, fifty. LLM serving is not M/M/1:
continuous batching changes service rate with batch composition, prompt/output
lengths are heavy-tailed, and memory is a second resource. The curve still teaches a
durable lesson: operating with no slack makes queueing extremely sensitive to small
load errors.

Memory amplifies the effect. If an admitted request holds on average $m$ bytes of KV,
Little's law suggests an average resident demand near $m\lambda T$ for a simplified
steady workload. As queueing increases $T$, more state remains live for longer. A
compute bottleneck can therefore turn into a memory crisis.

\section{Draw the state machine before choosing a policy}

A request moves through
\[
\textsc{New}\rightarrow\textsc{Queued}\rightarrow\textsc{Admitted}
\rightarrow\textsc{Running}\rightarrow\textsc{Terminal}\rightarrow\textsc{Retired}.
\]
Preemption adds
\[
\textsc{Running}\rightarrow\textsc{Preempted}\rightarrow\textsc{Queued}.
\]
Rejection is different:
\[
\textsc{New/Queued}\rightarrow\textsc{Rejected}.
\]
A rejected request never owns execution state. A preempted request has already
consumed service and may retain logical output even if its physical KV is reclaimed.

The distinction matters operationally. A service can report immediate overload for a
rejection. A preempted interactive request may exhibit a long mid-stream pause. A
timed-out request that remains physically resident is neither proper rejection nor
proper preemption; it is leaked work.

\EngineFigure{ch19-state-machine}{Admission, execution, preemption and retirement are distinct transitions. Physical KV may be reclaimed only through an ownership-safe path; rejection occurs before expensive execution state is established.}{fig:sched-state}

\section{Admission: promise, estimate, or gamble}

Suppose request $r$ has prompt length $S_r$, maximum output $O_r^{\max}$ and KV cost
$m_{kv}$ bytes/token. Worst-case reservation requires
\[
M_r^{\max}=(S_r+O_r^{\max})m_{kv}.
\]
Admit only if this fits and completion is memory-safe, but capacity can be wasted when
most outputs terminate early.

Optimistic admission reserves less. It may use current prompt KV plus a growth margin,
an output-length predictor, a percentile of historical output length, or no future
reservation at all. More requests start, but the system must have a recovery path
when predictions are wrong.

The useful design is not ``conservative versus optimistic.'' It is an explicit risk
budget. State what quantity is reserved, what uncertainty remains, what pressure
threshold triggers reclamation, and what user-visible outcome follows if the estimate
fails.

\section{Chunking makes admission easier to get wrong}

Chapter~\ref{ch:chunked-prefill} showed that a long request can enter through a small
first chunk. If admission considers only the next iteration, many long prompts can
simultaneously appear affordable. Their KV grows later and the allocator discovers
the debt after substantial compute has been spent.

Therefore separate
\[
B_{\rm step}\quad\text{from}\quad B_{\rm resident}.
\]
The first is an execution budget; the second is a capacity/admission budget. Current
vLLM exposes an option to reserve against full input length with chunked prefill to
reduce this form of over-admission. The durable rule is engine-independent: temporal
slicing does not erase future state growth.

\section{Pressure is a signal, not yet a victim}

Assume the allocator needs $\Delta b$ blocks for the next scheduled work and has only
$f<\Delta b$ free. The scheduler now has choices:
\begin{enumerate}
\item shrink or defer the work quantum;
\item evict unpinned prefix-cache entries;
\item preempt one or more active requests;
\item restore capacity from another tier or worker;
\item reject/not admit additional work;
\item fail if no legal recovery exists.
\end{enumerate}
These actions have different costs. Cached-but-inactive blocks are usually cheaper to
evict than active state. A request one token from completion is often a poor
recomputation victim. A low-priority batch request may be the correct victim even if
it arrived earlier.

\section{Preemption by recomputation}

Recomputation preemption preserves the request's logical token history, frees its
physical KV, and later reconstructs the KV by running the known tokens through the
model again. If the model execution is deterministic under the serving contract, the
reconstructed state represents the same causal history.

Let $C_r$ be the compute time required to rebuild the evicted prefix. A rough victim
cost is
\[
C_r \approx \frac{F_{\rm prefill}(S_r^{\rm committed})}{\pi_{\rm eff}},
\]
reduced by any prefix-cache portion that remains valid. Reclaim benefit is the number
of blocks actually made free, not the request's logical context length if blocks are
shared.

Recomputation spends accelerator compute but avoids host capacity and transfer
bandwidth. It is attractive when prefill is efficient and the victim's reconstructible
history is not enormous.

\section{Preemption by swapping}

Swap preemption moves KV to a lower memory tier and later restores it. For $M_r$ bytes
and effective bidirectional transfer bandwidth $\beta$, a first-order round-trip cost
is
\[
T_{\rm swap}\approx\frac{2M_r}{\beta}+T_{\rm setup}.
\]
This is derived and intentionally ignores overlap, topology contention, serialization
and page granularity.

Compare $T_{\rm swap}$ with recomputation time, but also price scarce host memory and
interconnect bandwidth. Long contexts with expensive recomputation may favor restore;
short contexts often favor recomputation. Chapters on KV storage and offloading treat
the tiering mechanisms in depth.

\section{A victim score should price damage, not age alone}

A conceptual victim score can be written
\begin{equation}
V_r =
\frac{R_r}
     {\epsilon+C_r+\alpha P_r+\beta U_r+\gamma D_r},
\label{eq:victim-score}
\end{equation}
where $R_r$ is reclaimable bytes, $C_r$ estimated restore/recompute cost, $P_r$ a
priority penalty, $U_r$ service already invested, and $D_r$ deadline urgency.
Coefficients encode service policy; $\epsilon$ prevents division by zero. Larger
$V_r$ means more capacity reclaimed per unit of damage.

Equation~\ref{eq:victim-score} is illustrative, not a production recommendation.
Simple policies such as latest-admitted or lowest-priority are easier to audit and can
be more robust than a poorly calibrated multi-factor score. The important step is to
name what the victim policy optimizes.

\section{Shared blocks make ``free this request'' nontrivial}

With prefix caching, request $A$ may reference 100 blocks but uniquely own only 20.
Preempting it cannot reclaim the 80 blocks still referenced by request $B$ or retained
by cache policy. Victim selection therefore needs \emph{reclaimable physical bytes},
not logical context bytes.

The Chapter~\ref{ch:paged-kv} refcount/generation contract supplies the answer.
Preemption releases the victim's references; only blocks whose reference/pin counts
reach the reclaimable state enter the free pool. A scheduler that predicts 100 blocks
of relief and receives 20 can immediately enter a preemption cascade.

\section{Avoid preemption cascades}

A cascade occurs when the scheduler repeatedly evicts requests but each victim frees
too little, or when newly admitted work consumes the reclaimed blocks before the
original pressure is resolved. This is a feedback failure.

Make reclamation transactional. Compute a target free-watermark $W$. Select a victim
set whose predicted unique reclaimable blocks reaches $W$, stop new admission while
reclamation is in progress, release references, observe actual free capacity, and
only then schedule the blocked allocation. If prediction was wrong, recompute the
plan rather than blindly continuing the original victim list.

\EngineFigure{ch19-pressure-loop}{Memory pressure enters a bounded recovery loop: evict cheap inactive cache first, select sufficient active victims only if necessary, observe actual reclaimed capacity, then resume admission.}{fig:sched-pressure-loop}

\section{Fairness is measured in service, not request count}

Two tenants can each have ten requests while one asks for 50 tokens/request and the
other 5,000. Round-robin by request is not equal service. Fair schedulers need a cost
unit.

Virtual Token Counter (VTC) scheduling accounts weighted input/output token service
per client and favors the backlogged client that has received less normalized service.
Its authors derive bounded service difference while retaining work conservation. The
important abstraction is a virtual service counter, not the exact formula: account
what scarce work each principal has received, then schedule against that accounting.

A production service may weight tenants by purchased capacity:
\[
\tilde s_i=\frac{s_i}{w_i},
\]
where $s_i$ is accumulated service and $w_i$ entitlement weight. Serving the smallest
$\tilde s_i$ approximates weighted fairness. Reset/decay policy matters; otherwise a
tenant that was idle for hours can accumulate an unrealistic credit advantage.

\section{Priority is not fairness}

Priority answers ``who should win when objectives conflict?'' Fairness answers ``how
should comparable principals share service over time?'' A strict priority scheduler
can starve low-priority work forever. A perfectly fair scheduler can violate an
emergency request's deadline.

Combine them explicitly. One approach chooses the highest eligible priority class,
then applies fair scheduling within that class, with aging that eventually promotes
old lower-priority work. Another reserves capacity shares by class. The policy must be
documented because it changes user-visible latency even though model correctness is
unchanged.

\section{Deadlines introduce value decay}

If a request's useful deadline is $d_r$ and predicted completion is later than
$d_r$, continuing it may have zero service value. Define slack
\[
\sigma_r=d_r-t_{\rm now}-\widehat T_r.
\]
Negative slack means the predictor expects a miss. Earliest-deadline-first can help
when work sizes are comparable, but a huge nearly-impossible request can monopolize
the server. Admission should therefore consider both deadline and estimated remaining
work.

Rejecting an impossible request early can increase goodput because it preserves
capacity for requests that can still meet their objectives.

\section{Goodput: count useful completions}

Raw throughput counts tokens regardless of latency. Goodput counts requests or tokens
that satisfy a declared service objective. For request indicator
\[
g_r =
\begin{cases}
1,&\mathrm{TTFT}_r\le\tau_{\rm TTFT}\ \land\
   \mathrm{TPOT}_r\le\tau_{\rm TPOT},\\
0,&\text{otherwise},
\end{cases}
\]
request goodput over interval $\Delta t$ is
\[
G=\frac{\sum_r g_r}{\Delta t}.
\]
The exact SLO may include percentile rather than per-request conditions; state it.

Goodput changes optimization behavior. Admitting one more request can increase raw
throughput while decreasing $G$ if it pushes many existing streams over their ITL
limits. Capacity planning therefore needs a safe operating region, not merely a peak
benchmark.

\section{SLO-aware admission as a feasibility test}

Before admission, estimate the marginal effect of request $r$ on memory and latency.
A conceptual test is
\[
\widehat M_{\rm active}+M_r^{\rm reserve}\le M_{\rm budget}
\]
and
\[
\widehat{\mathrm{TTFT}}_r\le\tau_{\rm TTFT},\qquad
\Delta\widehat{\mathrm{ITL}}_{\rm incumbents}\le\delta.
\]
If infeasible, defer within a bounded queue or reject immediately.

Predictions need uncertainty margins. A point estimate of 90 ms under a 100 ms SLO is
not safe when prediction error is $\pm40$ ms. Production admission can use calibrated
quantiles or conservative headroom. The cost of a false accept is wasted compute plus
damage to incumbents; the cost of a false reject is lost useful work.

\section{Bound the queue}

An unbounded queue converts overload into latency and memory. Even if queued requests
hold no KV, they hold client expectations and may expire before service. Define a
maximum queue length, maximum queue age, or both.

When rejecting, expose a retryable overload outcome and use randomized/exponential
backoff. Immediate synchronized retries can turn a small overload into a retry storm.
The server and client form one feedback system.

\section{Why feedback controllers oscillate}

Mooncake's production study gives a useful example. In a disaggregated prefill/decode
system, rejecting at the decode cluster after prefill wastes expensive work. Early
rejection at arrival is better, but a controller reacting to \emph{current} decode
load can oscillate because the request reaches decode only after prefill delay.

The authors instead predict decode load at the future handoff time. Their trace replay
of 23,000 requests at doubled speed reported fewer rejections with prediction-based
early rejection than with the baseline or instantaneous early-rejection policy.
Those are authors' measurements from their workload, not universal ratios.

The control lesson generalizes: if an action has delayed effect, regulate against the
state at the effect horizon, not only the state observed now.

\section{Cache-aware scheduling can conflict with fairness}

Chapter~\ref{ch:prefix-caching} made prefix locality valuable. SGLang exposes policies
including longest-prefix matching because warm-prefix work can reduce prefill cost.
But always serving the best cache hit can starve cold requests.

Treat locality as one scheduling feature, not the entire policy. A fair scheduler can
choose among eligible tenants first and then use prefix locality inside the chosen
tenant. Alternatively, add bounded locality credit that cannot overwhelm queue age.
This is an example of policy composition: optimize reuse subject to fairness, rather
than optimize reuse and hope fairness emerges.

\section{What current engines expose}

Current vLLM configuration documents FCFS and priority scheduling policy options and
contains modern machinery around deferred block release and scheduler accounting.
Historical vLLM versions used recomputation as the normal preemption path when KV
pressure made a running set infeasible. Because these internals evolve quickly, exact
current preemption claims must be pinned to source before this chapter is promoted.

Current SGLang source at snapshot
\texttt{7a719e9a65e7f52b012ab37211b5f174e5d3d38d} contains explicit scheduling-policy
logic, including cache-locality-oriented choices. This is useful evidence that
scheduling policy is a first-class runtime component, but source inspection does not
establish that every deployment uses the same policy.

\section{Failure tests for the scheduler}

A scheduler test suite needs adversarial traces, not only average throughput.
Test a tenant burst against a single interactive request; a long request against many
short requests; identical priority with different arrival order; cache-locality bias
against a cold request; a request whose output estimate is badly wrong; and a victim
whose blocks are mostly shared.

Then test invariants. No terminal request may regain execution ownership. No
preemption may free blocks still pinned elsewhere. Service counters must be monotone
under their accounting epoch. Queue age must not reset on preemption. Admission
failure must occur before expensive state allocation. A recomputed request must
produce the same model state/output under the backend's numerical contract.

\EngineFigure{sys-admission-flow}{Admission flowchart for a conceptual SLO-aware scheduler. An estimate is not a guarantee. Deferral must be bounded; rejecting work is different from silently timing it out.}{fig:sys-admission-flow}

\LedgerSection
\begin{ConstraintLedger}
Memory capacity & buys & Optimistic admission plus bounded reclamation can run more requests than worst-case reservation while recovering from prediction error.\\
Compute & spends & Recomputation destroys resident KV and later pays prefill again; cascades can turn overload into thrashing.\\
Memory bandwidth & spends & Swap preemption moves KV across memory tiers and can contend with inference traffic.\\
Latency & moves & Rejection protects incumbents; preemption transfers delay to victims; fairness and priority decide distribution.\\
Correctness & risks & Shared-block reclamation, stale ownership, premature retirement or nondeterministic reconstruction can corrupt state.\\
Cost & trades & SLO-aware admission can improve useful completions per accelerator-hour; overly conservative rejection strands capacity.\\
\end{ConstraintLedger}

\LabSection{a finite-capacity scheduler that can prove whom it harmed}

The Rust lab in \texttt{code/labs/ch19-scheduling/} models two tenants, finite KV
blocks, weighted service counters, bounded admission and recomputation preemption.
Requests declare prompt blocks, remaining steps, priority and optional deadline. The
simulator records admissions, service, preemptions, rejections and completions.

Its oracle is deliberately independent of policy: conservation checks that each
completed request receives exactly its declared service units, and ownership checks
that allocated blocks never exceed capacity. A deterministic burst demonstrates that
rejecting one infeasible arrival can preserve completion of incumbents.

Broken-policy exercises are concrete: remove tenant normalization and the fairness
test fails; free a victim's logical blocks instead of unique blocks and capacity
accounting fails; reset queue age after preemption and the starvation test fails; let
admission exceed capacity and the safety invariant fails.

\HappenedSection{overload became an explicit product policy}

Early serving stacks often treated OOM as an exceptional condition. Modern engines
turn resource pressure into scheduler state: block pools report capacity, queues are
ordered by explicit policies, requests can be deferred or preempted, and production
systems use admission to protect latency objectives.

The measured llama.cpp case remains valuable because it shows the opposite behavior
in one pinned server revision: cache exhaustion propagated to every active slot.
That is not a statement about all current llama.cpp configurations; it is a
reproducible historical example of why failure policy belongs in the architecture.

Mooncake adds the distributed version of the lesson. Rejecting only after prefill
wastes upstream work; rejecting based on instantaneous downstream load can oscillate;
prediction must account for the delay between decision and effect.

\FrontierSection{2026-10-05}

Scheduling research is shifting from maximizing token throughput toward multi-resource
goodput: KV capacity, prefill/decode compute, cache locality, tenant fairness and SLO
deadlines interact. VTC provides a formal token-service fairness model. DistServe and
Mooncake make SLO-aware admission central in disaggregated serving. Current SGLang
keeps scheduling policy explicit in source, while current vLLM exposes FCFS/priority
policy and increasingly detailed KV/scheduler accounting.

The frontier question is not one universal policy. It is whether a runtime can expose
enough cost and ownership information to compose policies safely: locality subject to
fairness, priority subject to starvation bounds, optimistic admission subject to a
recovery budget, and throughput subject to latency goodput.

\ExercisesSection
\begin{enumerate}
\item \emph{Derivation.} From equation~\eqref{eq:sched-mm1}, solve for utilization
when mean system time is 5, 20 and 50 service times.
\item \emph{Trace.} Four requests each need 200 prompt cells and up to 1,200 output
cells. Compute worst-case total capacity. With 4,096 cells and equal progress, at
approximately how many generated tokens/request is capacity exhausted?
\item \emph{Victim selection.} Requests A and B reference 100 blocks each. A uniquely
owns 20 and costs 5 ms to reconstruct; B uniquely owns 70 and costs 40 ms. Compare
reclaimable-blocks-per-ms and explain what the ratio omits.
\item \emph{Implementation.} Extend the lab with shared block refcounts. Acceptance:
preemption frees only unique/unpinned blocks; actual reclaimed capacity is observed
before new admission; no preemption cascade can exceed a configured victim bound.
\item \emph{Fairness.} Construct a short trace where weighted fair service looks
unfair over three iterations but converges over a longer backlogged interval.
\item \emph{Falsification.} Construct a workload where strict SLO-aware rejection
reduces goodput relative to optimistic admission. Identify which predictor error or
burst structure makes the conservative policy lose.
\end{enumerate}

\Needspace{18\baselineskip}
\section{Worked problems}

\Needspace{12\baselineskip}
\subsection{Queueing ratios}
\textbf{Problem.} \emph{Derivation.} From equation~\eqref{eq:sched-mm1}, solve for utilization
when mean system time is 5, 20 and 50 service times.

\textbf{Solution.} From $T/(1/\mu)=1/(1-\rho)=k$, $\rho=1-1/k$. For $k=5,20,50$ the
utilizations are 0.80, 0.95 and 0.98.

\Needspace{12\baselineskip}
\subsection{Opening capacity}
\textbf{Problem.} \emph{Trace.} Four requests each need 200 prompt cells and up to 1,200 output
cells. Compute worst-case total capacity. With 4,096 cells and equal progress, at
approximately how many generated tokens/request is capacity exhausted?

\textbf{Solution.} Worst-case cells are $4(200+1200)=5600$. With 4,096 cells and four 200-token
prompts, 3,296 cells remain for outputs, or 824 output cells/request under equal
progress. This reproduces the pressure point in the recorded setup.

\Needspace{12\baselineskip}
\subsection{Victim efficiency}
\textbf{Problem.} \emph{Victim selection.} Requests A and B reference 100 blocks each. A uniquely
owns 20 and costs 5 ms to reconstruct; B uniquely owns 70 and costs 40 ms. Compare
reclaimable-blocks-per-ms and explain what the ratio omits.

\textbf{Solution.} A reclaims 20/5=4 unique blocks/ms of reconstruction cost. B reclaims
70/40=1.75 blocks/ms, so this narrow metric chooses A. It omits priority, deadline,
service already invested, cache reuse, transfer alternatives and whether the freed
blocks are sufficient to satisfy the blocked allocation.

\Needspace{12\baselineskip}
\subsection{Conservative-policy falsification}
\textbf{Problem.} \emph{Falsification.} Construct a workload where strict SLO-aware rejection
reduces goodput relative to optimistic admission. Identify which predictor error or
burst structure makes the conservative policy lose.

\textbf{Solution.} If the predictor systematically overestimates output length, strict reservation can
reject short requests that would have completed safely. Optimistic admission can then
produce higher SLO-qualified completions when actual outputs terminate early.

